% Standalone chunk A.2: Concentration chain (POST-FIX, Round 2).
% Covers: Preliminaries, F2/F3, Lemma K_inverse, Lemma quad_conf,
%         Lemma G_unbiased_conf (scaled-identity bias, exact),
%         Rem bias_floor_sharp,
%         Lemma G_bound_conf (observable cA_t event),
%         Theorem matrix_bernstein_conf (predictable truncation MDS),
%         Corollary projector_conf (with m_k >= log(2d/delta) hypothesis),
%         Proposition segment_factorization_conf.
% Source: conf.tex lines ~1057-1402 (post-fix).
% Shared preamble for standalone chunks. Do not compile directly; \input{} it.
\documentclass{article}
\usepackage[utf8]{inputenc}
\usepackage[T1]{fontenc}
\usepackage{lmodern}
\usepackage[margin=1in]{geometry}
\usepackage{amsmath,amssymb,amsfonts,amsthm,bm,mathtools}
\usepackage{enumitem}
\usepackage{microtype}
\usepackage{xcolor}
\usepackage{natbib}
\usepackage{booktabs}
\usepackage{hyperref}

\newcommand{\R}{\mathbb{R}}
\newcommand{\N}{\mathbb{N}}
\newcommand{\E}{\mathbb{E}}
\newcommand{\Prob}{\mathbb{P}}
\newcommand{\op}{\mathrm{op}}
\DeclareMathOperator{\sign}{sign}
\newcommand{\cA}{\mathcal{A}}
\newcommand{\cF}{\mathcal{F}}
\newcommand{\cH}{\mathcal{H}}
\newcommand{\cT}{\mathcal{T}}
\newcommand{\cI}{\mathcal{I}}
\newcommand{\cW}{\mathcal{W}}
\newcommand{\cTprobe}{\mathcal{T}_{\mathrm{probe}}}
\newcommand{\cE}{\mathcal{E}}
\newcommand{\norm}[1]{\left\lVert #1 \right\rVert}
\newcommand{\tilO}{\widetilde{\mathcal{O}}}
\newcommand{\DynReg}{\mathrm{DynReg}}
\newcommand{\rank}{\mathrm{rank}}
\newcommand{\range}{\mathrm{range}}
\DeclareMathOperator{\TopEig}{TopEig}
\DeclareMathOperator*{\argmax}{arg\,max}
\DeclareMathOperator{\tr}{tr}

\newtheorem{assumption}{Assumption}
\newtheorem{theorem}{Theorem}
\newtheorem{corollary}{Corollary}
\newtheorem{proposition}{Proposition}
\newtheorem{lemma}{Lemma}
\newtheorem{remark}{Remark}


\title{Chunk A.2: Concentration of the lifted probe moment\\(post-fix, Round 4 verification)}
\author{(excerpt for standalone review)}
\date{}

\begin{document}
\maketitle

\textbf{Setup context (not in chunk).} The learner plays $x_t \in \cA_t
\subset \R^d$ with $\|x\| \le R_\cA$ and observes
$y_t = x_t^\top\theta_t + \varepsilon_t$, $\theta_t = B_k^\star w_t$,
$B_k^\star \in \R^{d\times r}$ orthonormal, $\|w_t\|\le S_w$,
$\varepsilon_t$ sub-Gaussian with $\psi_2$-norm $\sigma_\varepsilon$
(variance $\sigma_\varepsilon^2$, known up to a plugin $\hat\sigma^2$;
misspec $\delta_\sigma := \hat\sigma^2-\sigma_\varepsilon^2$), and
$\|\E[\varepsilon_t\theta_t\mid\cH_{t-1}]\|\le\epsilon_\times$
(state--noise coupling). Probes are iid $u_t\sim\mathcal N(0,I_d)$,
truncated at $\|u\|\le L$ (handled in Appendix~C chunk). The probe
statistic is $s_t := y_t^2 - \hat\sigma^2$ and the lifted sample
$G_t := \mathcal K^{-1}(s_t u_tu_t^\top)$, where $\mathcal K:M\mapsto\E[(u^\top M u)uu^\top]$.

\textbf{What changed since Round~3 (R4 delta, this chunk).}
\begin{itemize}[leftmargin=*,itemsep=1pt,topsep=1pt]
\item \textbf{Cor.~\ref{cor:projector_conf}} hypotheses rewritten as three
  explicit itemized conditions (Freedman domination, Large-sample,
  Probe-bias smallness), replacing R3's parenthetical ``provided
  $b_\sigma\le\lambda_{\min}/4$''. The probe-bias smallness now
  explicitly includes $\|\Theta_k\|_\op$, not just $b_\sigma$.
\item Confidence-chain wording corrected: the Freedman event itself
  holds w.p.\ $\ge 1-2\delta$ (Freedman $\delta$ + truncation-event $\delta$
  from Lem.~\ref{lem:G_bound_conf}), not $\ge 1-\delta$ as R3's text implied.
\end{itemize}

\textbf{What changed since Round~2 (R3 delta, retained).}
\begin{itemize}[leftmargin=*,itemsep=1pt,topsep=1pt]
\item \textbf{Cor.~\ref{cor:projector_conf}} restated: the R2 hypothesis
``assume $\|\widehat M_k-\bar M_k^{\mathrm{probe}}\|_\op\le\lambda_{\min}/2$''
was a random event used as a deterministic assumption. It is replaced
by a proper deterministic sample-size condition
$m_k\ge(16R_X/\lambda_{\min})^2\log(2d/\delta)$ that \emph{implies}
the event w.h.p.\ via the Freedman bound.
\item \textbf{Cor.~\ref{cor:projector_conf}} bound augmented with an
explicit $4\|\Theta_k\|_\op/\lambda_{\min}$ term; R2 absorbed this
silently into the $C_{\mathrm{sub}}$ inflation, which made downstream
references (e.g.\ Cor.~rank\_adaptive, Round~2 chunk~A.5) unclear.
\end{itemize}

\textbf{What changed since Round~1 (delta highlights, retained from R2).}
\begin{itemize}[leftmargin=*,itemsep=1pt,topsep=1pt]
\item \textbf{Lem.~\ref{lem:G_unbiased_conf}}: bias refined from conservative
$\|\widetilde B\|_\op\le|\delta_\sigma|L^2+2L^3\epsilon_\times$ to
$\widetilde B_t=-\delta_\sigma/(d+2)\cdot I_d$ (scaled identity, exact),
using Gaussian odd-moment vanishing for the $\epsilon_\times$ cross term.
\item \textbf{Lem.~\ref{lem:G_bound_conf}}: truncation event rewritten as
the \emph{observable} $\cA_t:=\{\|u_t\|\le L\}\cap\{|y_t|\le L_y\}$
(both $\sigma(u_t,y_t)$-measurable); probabilities accounted with
$L^2=2d\log(4T/\delta)$, $L_\varepsilon=\sigma_\varepsilon\sqrt{2\log(4T/\delta)}$.
\item \textbf{Thm.~\ref{thm:matrix_bernstein_conf}}: Freedman applied to the
\emph{predictably-truncated MDS} $\tilde X_t:=G_t\mathbf 1\{\cA_t\}-\E[G_t\mathbf 1\{\cA_t\}\mid\cH_{t-1}]$;
a deterministic offset $\Theta_k$ is introduced to transfer the bound
back to $\widehat M_k$ on the high-probability event $\cE:=\bigcap_t\cA_t$.
\item \textbf{Cor.~\ref{cor:projector_conf}}: hypothesis strengthened
with $m_k\ge\log(2d/\delta)$ so that the linear Freedman term is dominated
by the square-root term; truncation bias $\|\Theta_k\|_\op$ accounted
explicitly (bounded by $O(d\sqrt{\delta/T})$ in the Appendix~C chunk).
\end{itemize}

\section*{Preliminaries (excerpt)}

$\cI_k = [\tau_{k-1}, \tau_k)$ denotes segment $k$ of length
$\ell_k := |\cI_k|$, $\sum_k \ell_k = T$. Inside $\cI_k$ the learner
performs $m_k = |\cT_k|$ probe rounds ($\cT_k := \cTprobe\cap\cI_k$)
and $n_k = \ell_k - m_k$ exploitation rounds ($E_k := \cI_k\setminus\cT_k$).
Write $\widehat P_t = \widehat U_t \widehat U_t^\top$ and
$P_k^\star = B_k^\star B_k^{\star\top}$ for the estimated and true rank-$r$
projectors, and $\varepsilon_k := \|\widehat P_k - P_k^\star\|_{\op}$.
The \emph{predictable second moment} is
$\widetilde M_t := \E[\theta_t\theta_t^\top\mid\cH_{t-1}]
= B_k^\star \E[w_tw_t^\top\mid\cH_{t-1}](B_k^\star)^\top$, and the
probe-time average is
$\bar M_k^{\mathrm{probe}} := m_k^{-1}\sum_{t\in\cT_k}\widetilde M_t$.
By the stable-LDS assumption and $\Sigma_\eta\succ 0$,
$\bar M_k^{\mathrm{probe}}$ has rank $r$ with $r$-th eigenvalue
$\ge\lambda_{\min}>0$ uniformly in $k$.

\begin{lemma}[Predictable probe-time excitation]
\label{lem:probe_excitation_conf}
For each segment $k$,
$\lambda_r(m_k^{-1}\sum_{t\in\cT_k}\E[w_tw_t^\top\mid\cH_{t-1}])
\ge \lambda_{\min}>0$ a.s.
\end{lemma}
\begin{proof}
By stability and $\Sigma_\eta\succ 0$, the Lyapunov equation
$\Sigma_k = A_k\Sigma_k A_k^\top + \Sigma_{\eta,k}$ has unique
positive-definite solution with $\lambda_r(\Sigma_k)\ge
\lambda_r(\Sigma_{\eta,k})/(1-(1-\alpha_0)^2)=:\underline\lambda$.
Since $\E[w_tw_t^\top\mid\cH_{t-1}]\to\Sigma_k$ exponentially inside
$\cI_k$, the probe-time average satisfies
$\lambda_r(m_k^{-1}\sum_{t\in\cT_k}\E[w_tw_t^\top\mid\cH_{t-1}])\ge\underline\lambda/2=:\lambda_{\min}$
past an $O(1)$ warm-up.
\end{proof}

\noindent\textbf{F2 (Windowed self-normalised bound).}\ On projected
covariates $z_t = \widehat U_{t-1}^\top x_t\in\R^r$, the windowed ridge
estimator $\widehat a_t = \widetilde V_t^{-1}\widetilde b_t$ satisfies,
with probability $\ge 1-\delta$ and simultaneously for all $t$,
$\|\widehat a_t - a_t^\star\|_{\widetilde V_t}\le\beta^{(r,W)}_t$
with $\beta^{(r,W)}_t = \sigma_\varepsilon\sqrt{r\log(1+WR_\cA^2/(\lambda r)) + 2\log(1/\delta)} + \sqrt\lambda S_w$
(projected Abbasi-Yadkori with sliding-window correction of Russac et al.).

\noindent\textbf{F3 (Elliptical potential).}\ For any adapted
$z_1,\dots,z_T\in\R^r$ with $\|z_t\|\le R_\cA$ and
$\widetilde V_t = \lambda I + \sum_{s\in\cW_t} z_s z_s^\top$,
$\sum_{t=1}^T \|z_t\|^2_{\widetilde V_t^{-1}}
\le 2r\log(1 + W R_\cA^2/(\lambda r)) =: 2r L_W$.

\section*{Concentration chain}

\begin{lemma}[Closed-form inverse of $\mathcal K$]
\label{lem:K_inverse}
For isotropic Gaussian probes $u\sim\mathcal N(0,I_d)$,
$\mathcal K(M) = 2M + \tr(M)\,I_d$ and
$\mathcal K^{-1}(N) = \tfrac12 N - \tfrac{\tr(N)}{2(d+2)}\,I_d$, so
$\|\mathcal K^{-1}\|_\op \le 1$.
\end{lemma}
\begin{proof}
Expand $u^\top M u = \sum_{ij} M_{ij} u_i u_j$ and use Isserlis
$\E[u_i u_j u_k u_l] = \delta_{ij}\delta_{kl} +
\delta_{ik}\delta_{jl} + \delta_{il}\delta_{jk}$ to compute
$\E[(u^\top M u)u_k u_l] = M_{kl}+M_{lk}+\tr(M)\delta_{kl} = 2M_{kl}
+ \tr(M)\delta_{kl}$ (symmetric $M$). Inverting
$N = 2M + \tr(M)I_d$: $\tr(N) = (d+2)\tr(M)$, hence
$M = \tfrac12 N - \tfrac{\tr(N)}{2(d+2)}I_d$. The operator norm
follows from $|\tr(N)|/(d+2)\le\|N\|_\op$ for symmetric $N$.
\end{proof}

\begin{lemma}[Quadratic measurement identity]
\label{lem:quad_conf}
On every probe round $t\in\cTprobe$, with
$\delta_\sigma := \hat\sigma^2 - \sigma_\varepsilon^2$,
\[
\E[s_t\mid\sigma(\cH_{t-1},u_t)]
= u_t^\top \widetilde M_t u_t + 2 u_t^\top m_t - \delta_\sigma,\qquad \|m_t\|\le\epsilon_\times,
\]
where $m_t:=\E[\varepsilon_t\theta_t\mid\cH_{t-1},u_t]$. Under exact
probe conditions ($\delta_\sigma=\epsilon_\times=0$) this collapses to
$\E[s_t\mid\cdot] = u_t^\top\widetilde M_t u_t$.
\end{lemma}
\begin{proof}
Expand $y_t^2 = (u_t^\top\theta_t)^2 + 2\varepsilon_t u_t^\top\theta_t
+ \varepsilon_t^2$ and condition on $\sigma(\cH_{t-1},u_t)$. Middle
term contributes $2u_t^\top m_t$ with $\|m_t\|\le\epsilon_\times$ by
assumption; last term contributes $\sigma_\varepsilon^2$, and
subtracting $\hat\sigma^2$ leaves $-\delta_\sigma$. Independence of
$u_t$ from $(\theta_t,\cH_{t-1})$ gives the first term.
\end{proof}

\begin{lemma}[Lifted probe sample is near-unbiased --- scaled-identity bias]
\label{lem:G_unbiased_conf}
Let $G_t := \mathcal K^{-1}(s_t u_tu_t^\top)$. Under $u_t\sim\mathcal N(0,I_d)$
(untruncated; truncation handled separately),
$\E[G_t\mid\cH_{t-1}] = \widetilde M_t + \widetilde B_t$ where
\[
\widetilde B_t \;=\; -\frac{\delta_\sigma}{d+2}\,I_d \qquad(\text{scaled identity, exact}),
\]
so $\|\widetilde B_t\|_\op = |\delta_\sigma|/(d+2)$. In particular
$\widetilde B_t$ shifts all eigenvalues of $\E[G_t\mid\cH_{t-1}]$ by
the same amount and does \emph{not} rotate eigenvectors.
\end{lemma}
\begin{proof}
Linearity of $\mathcal K^{-1}$ and Lem.~\ref{lem:quad_conf} give
\[
\E[G_t\mid\cH_{t-1}]
= \mathcal K^{-1}(\mathcal K(\widetilde M_t))
+ 2\,\mathcal K^{-1}(\E[(u_t^\top m_t)u_tu_t^\top\mid\cH_{t-1}])
- \delta_\sigma\,\mathcal K^{-1}(\E[u_tu_t^\top]).
\]
The first term equals $\widetilde M_t$ (Lem.~\ref{lem:K_inverse}).
The cross term vanishes: for any $\cH_{t-1}$-measurable $m_t$,
$u_t\sim\mathcal N(0,I_d)$ has all third moments zero by symmetry, so
$\E[(u_t^\top m_t)u_{t,k}u_{t,l}]=\sum_i m_{t,i}\,\E[u_{t,i}u_{t,k}u_{t,l}]=0$
identically. Hence the $\epsilon_\times$-dependent term contributes
zero to $\widetilde B_t$. For the third term, $\E[u_tu_t^\top]=I_d$,
so Lem.~\ref{lem:K_inverse} gives
$-\delta_\sigma\mathcal K^{-1}(I_d) = -\delta_\sigma[I_d/2 - d/(2(d+2))\,I_d]
= -\delta_\sigma/(d+2)\cdot I_d$. Combining yields the scaled-identity form.
\end{proof}

\begin{remark}[Sharpened bias floor]
\label{rem:bias_floor_sharp}
Under isotropic Gaussian probes the state--noise coupling
$\epsilon_\times$ contributes \emph{zero} bias to $\widetilde B_t$ (by
Gaussian odd-moment vanishing), in contrast to the previously-stated
conservative bound $|\delta_\sigma|L^2 + 2L^3\epsilon_\times$. The
conservative bound remains valid for non-Gaussian probes lacking the
odd-moment-vanishing property, and is retained in
Cor.~\ref{cor:projector_conf} for generality; for Gaussian probes
specifically, $\widetilde B_t$ is exactly a scaled identity and
$\Delta_\sigma$ reduces to the sharper bound
$\Delta_\sigma^{\mathrm{sharp}} = 4|\delta_\sigma|/((d+2)\lambda_{\min})$.
\end{remark}

\begin{lemma}[Observable a.s.\ bound on $G_t$]
\label{lem:G_bound_conf}
Define the \emph{observable} truncation event
$\cA_t := \{\|u_t\|_2\le L\}\cap\{|y_t|\le L_y\}$ with
$L_y := L\,S_w + L_\varepsilon$ and
$L_\varepsilon := \sigma_\varepsilon\sqrt{2\log(4T/\delta)}$.
Both indicators $\mathbf 1\{\|u_t\|\le L\}$, $\mathbf 1\{|y_t|\le L_y\}$
are $\sigma(u_t,y_t)$-measurable, hence observable at round $t$. On
$\cA_t$,
\[
  |s_t|\le R_s := L_y^2 + \hat\sigma^2, \qquad
  \|G_t\|_\op \le R_X := L^2 R_s + S_w^2.
\]
Moreover, choosing $L^2=2d\log(4T/\delta)$ and
$L_\varepsilon=\sigma_\varepsilon\sqrt{2\log(4T/\delta)}$ gives
$\Pr(\cA_t^c\mid\cH_{t-1})\le\delta/T$, hence
$\Pr(\bigcap_{t\in\cTprobe}\cA_t)\ge 1-\delta$.
\end{lemma}
\begin{proof}
Since $\theta_t=B_k^\star w_t$ with $B_k^\star$ orthonormal and
$\|w_t\|\le S_w$, $\|\theta_t\|_2\le S_w$. On $\cA_t$,
$|y_t|\le L_y$ so $|s_t|=|y_t^2-\hat\sigma^2|\le y_t^2+\hat\sigma^2\le R_s$;
using $\|\mathcal K^{-1}\|_\op\le 1$ and $\|u_tu_t^\top\|_\op\le L^2$,
$\|G_t\|_\op\le|s_t|L^2+S_w^2=R_sL^2+S_w^2=R_X$ (the $S_w^2$
absorbs the centering term $\|\widetilde M_t\|_\op\le S_w^2$).

For the probability bound, $\chi^2_d$ tail with $L^2=2d\log(4T/\delta)$
gives $\Pr(\|u_t\|>L)\le\delta/(4T)$. On $\{\|u_t\|\le L\}$,
$|y_t|\le L S_w+|\varepsilon_t|$, so
$\{|y_t|>L_y\}\cap\{\|u_t\|\le L\}\subseteq\{|\varepsilon_t|>L_\varepsilon\}$,
and sub-Gaussian tail gives $\Pr(|\varepsilon_t|>L_\varepsilon)\le\delta/(2T)$.
Union: $\Pr(\cA_t^c)\le\delta/(4T)+\delta/(2T)\le\delta/T$. Union over
$|\cTprobe|\le T$ gives $\Pr(\bigcap\cA_t)\ge 1-\delta$.
\end{proof}

\begin{theorem}[Matrix Freedman for $\widehat M_k$ via predictable truncation]
\label{thm:matrix_bernstein_conf}
Define the predictably-truncated lifted sample
$\tilde G_t := G_t\cdot\mathbf 1\{\cA_t\}$ and
$\tilde X_t := \tilde G_t - \E[\tilde G_t\mid\cH_{t-1}]$.
Then $\{\tilde X_t\}_{t\in\cT_k}$ is an MDS with
$\|\tilde X_t\|_\op\le 2R_X$ a.s.\ and
$\|\E[\tilde X_t^2\mid\cH_{t-1}]\|_\op\le R_X^2$ a.s. With probability
$\ge 1-\delta$,
\begin{equation}
  \Bigl\|m_k^{-1}\textstyle\sum_{t\in\cT_k}\tilde X_t\Bigr\|_\op
  \;\le\; 2R_X\sqrt{\log(2d/\delta)/m_k} + \tfrac{2R_X\log(2d/\delta)}{3m_k}.
\end{equation}
Transferring to the algorithm's estimator
$\widehat M_k=m_k^{-1}\sum_{t\in\cT_k}G_t$: on the high-probability
event $\cE := \bigcap_{t\in\cT_k}\cA_t$ (mass $\ge 1-\delta$ by
Lem.~\ref{lem:G_bound_conf}), $\tilde G_t = G_t$ for every $t\in\cT_k$,
so Freedman applies to $\widehat M_k$ up to a deterministic offset
$\Theta_k := m_k^{-1}\sum_{t\in\cT_k}\E[G_t\mathbf 1\{\cA_t^c\}\mid\cH_{t-1}]$
of operator norm $O(d\sqrt{\delta/T})$ (Appendix~C chunk). With
probability $\ge 1-2\delta$,
\begin{equation}
\|\widehat M_k - \bar M_k^{\mathrm{probe}} - \widetilde B\|_\op
\le 2R_X\sqrt{\log(2d/\delta)/m_k} + \frac{2R_X\log(2d/\delta)}{3m_k}
+ \|\Theta_k\|_\op.
\label{eq:matrix_bernstein}
\end{equation}
\end{theorem}
\begin{proof}
\emph{MDS property.} $\mathbf 1\{\cA_t\}$ is $\sigma(u_t,y_t)$-measurable
hence $\cH_t$-measurable, so $\tilde G_t$ is $\cH_t$-measurable and
$\E[\tilde X_t\mid\cH_{t-1}]=0$.

\emph{A.s.\ bound.} On $\cA_t$, $\|\tilde G_t\|_\op=\|G_t\|_\op\le R_X$
by Lem.~\ref{lem:G_bound_conf}. On $\cA_t^c$, $\tilde G_t=0$, so
$\|\tilde G_t\|_\op\le R_X$ unconditionally. The centered
$\|\tilde X_t\|_\op\le\|\tilde G_t\|_\op+\|\E[\tilde G_t\mid\cH_{t-1}]\|_\op\le 2R_X$.

\emph{Predictable variance.}
$\|\E[\tilde X_t^2\mid\cH_{t-1}]\|_\op\le\|\E[\tilde G_t^2\mid\cH_{t-1}]\|_\op\le R_X^2$
(follows from $\tilde G_t^2\preceq R_X^2 I$ on $\cA_t$, and zero on $\cA_t^c$).

\emph{Freedman.} Matrix Freedman (Tropp 2011) with the above bounds
gives the first displayed inequality on a $1-\delta$ event.

\emph{Transfer.} Let $\bar{\tilde G}_k:=m_k^{-1}\sum\tilde G_t$,
$\bar\mu_k:=m_k^{-1}\sum\E[\tilde G_t\mid\cH_{t-1}]$. On $\cE$,
$\bar{\tilde G}_k=\widehat M_k$ and
$\bar\mu_k = m_k^{-1}\sum\E[G_t\mid\cH_{t-1}]-\Theta_k
= \bar M_k^{\mathrm{probe}}+\widetilde B-\Theta_k$ (using
Lem.~\ref{lem:G_unbiased_conf}). Rearranging,
$\widehat M_k-\bar M_k^{\mathrm{probe}}-\widetilde B
= m_k^{-1}\sum\tilde X_t - \Theta_k$ on $\cE$, and triangle inequality
yields \eqref{eq:matrix_bernstein} on $\cE\cap\{\text{Freedman}\}$,
mass $\ge 1-2\delta$.
\end{proof}

\begin{corollary}[Projector concentration]
\label{cor:projector_conf}
Let $b_\sigma := |\delta_\sigma|L^2 + 2L^3\epsilon_\times$ denote the
(deterministic) conservative upper bound on $\|\widetilde B\|_\op$ valid
for any probe distribution (Lem.~\ref{lem:G_unbiased_conf}).
Assume the three conditions:
\begin{itemize}[leftmargin=1.5em,itemsep=1pt,topsep=1pt]
\item (\emph{Freedman domination})\ $m_k\ge\log(2d/\delta)$, so the
  linear Freedman term is at most the square-root term;
\item (\emph{Large-sample})\ $m_k\ge(16R_X/\lambda_{\min})^2\log(2d/\delta)$,
  so the Freedman square-root bound is $\le\lambda_{\min}/4$ (w.p.\
  $\ge 1-2\delta$);
\item (\emph{Probe-bias smallness})\ $b_\sigma+\|\Theta_k\|_\op\le\lambda_{\min}/4$,
  so $\|\widetilde B\|_\op+\|\Theta_k\|_\op\le\lambda_{\min}/4$, giving
  $\|\widehat M_k-\bar M_k^{\mathrm{probe}}\|_\op\le\lambda_{\min}/2$ at
  the same confidence via \eqref{eq:matrix_bernstein}.
\end{itemize}
Then with probability $\ge 1-2\delta$,
\begin{equation}
\varepsilon_k := \|\widehat P_k - P_k^\star\|_\op
\le C_{\mathrm{sub}}\sqrt{\log(2d/\delta)/m_k} + \Delta_\sigma + \frac{4\|\Theta_k\|_\op}{\lambda_{\min}},\
C_{\mathrm{sub}} := \frac{8 R_X}{\lambda_{\min}},\
\Delta_\sigma := \frac{4(|\delta_\sigma|L^2+2L^3\epsilon_\times)}{\lambda_{\min}}.
\label{eq:proj_bound_conf}
\end{equation}
The truncation-bias contribution $4\|\Theta_k\|_\op/\lambda_{\min} = O(d\sqrt{\delta/T}/\lambda_{\min})=o(1)$
for $\delta$ polynomial in $1/T$ (Prop.~theta\_bound, App.~C chunk). It is kept explicit here so
downstream corollaries (e.g.\ Cor.~rank\_adaptive) can reference it.
For Gaussian probes, $\Delta_\sigma$ specializes to the sharper
$\Delta_\sigma^{\mathrm{sharp}}=4|\delta_\sigma|/((d+2)\lambda_{\min})$
(Lem.~\ref{lem:G_unbiased_conf}, Rem.~\ref{rem:bias_floor_sharp}).
\end{corollary}
\begin{proof}
By Lem.~\ref{lem:G_unbiased_conf},
$\|\widetilde B\|_\op\le b_\sigma:=|\delta_\sigma|L^2+2L^3\epsilon_\times$
(conservative, any probe distribution). Combining with
Thm.~\ref{thm:matrix_bernstein_conf}, with probability $\ge 1-2\delta$,
\[
\|\widehat M_k - \bar M_k^{\mathrm{probe}}\|_\op
\le 2R_X\sqrt{\log(2d/\delta)/m_k}
+ \frac{2R_X\log(2d/\delta)}{3m_k}
+ \|\Theta_k\|_\op + b_\sigma.
\]
Under $m_k\ge\log(2d/\delta)$, $2R_X\log(2d/\delta)/(3m_k)\le 2R_X\sqrt{\log(2d/\delta)/m_k}/3$,
absorbing into a constant inflation of $C_{\mathrm{sub}}$.
$\|\Theta_k\|_\op = O(d\sqrt{\delta/T})$ is polynomially small
(Prop.~theta\_bound in Appendix~C chunk). $\bar M_k^{\mathrm{probe}}$
has rank exactly $r$ with $r$-th eigenvalue $\ge\lambda_{\min}$
(Prop.~\ref{prop:segment_factorization_conf}); Davis--Kahan gives
$\|\widehat P_k-P_k^\star\|_\op\le 4\|\widehat M_k-\bar M_k^{\mathrm{probe}}\|_\op/\lambda_{\min}$,
yielding \eqref{eq:proj_bound_conf}.
\end{proof}

\begin{proposition}[Segment factorisation]
\label{prop:segment_factorization_conf}
$\bar M_k^{\mathrm{probe}} = B_k^\star\bar S_k^{\mathrm{probe}}(B_k^\star)^\top$
with $\bar S_k^{\mathrm{probe}}:=m_k^{-1}\sum_{t\in\cT_k}\E[w_tw_t^\top\mid\cH_{t-1}]$.
In particular $\range(\bar M_k^{\mathrm{probe}}) = \range(B_k^\star)$.
\end{proposition}
\begin{proof}
$\widetilde M_t = B_k^\star\E[w_tw_t^\top\mid\cH_{t-1}](B_k^\star)^\top$
for $t\in\cI_k$; average over $\cT_k$ and factor $B_k^\star$. Range
equality uses $\lambda_r(\bar S_k^{\mathrm{probe}})\ge\lambda_{\min}>0$
(Lem.~\ref{lem:probe_excitation_conf}).
\end{proof}

This chain---Lemmas~\ref{lem:K_inverse}--\ref{lem:G_bound_conf},
Theorem~\ref{thm:matrix_bernstein_conf}, Corollary~\ref{cor:projector_conf},
Proposition~\ref{prop:segment_factorization_conf}---gives a
self-contained upper bound on subspace recovery at rate
$O(1/\sqrt{m_k})+\Delta_\sigma$, matched (up to $\sqrt{\log d}$) by
the minimax lower bound (chunk A.3).

\end{document}
